\section{Related Work}
\label{sec:related}

\paragraph{Regression verification and equivalence checking}
\citet{godlin2009regression} introduced \emph{regression verification}: proving that two
closely related programs produce the same output for every input. They do it by pairing
recursive calls and checking loop-free bodies with a bounded model checker.
SymDiff~\citep{lahiri2012symdiff} translates both versions to an intermediate verification
language and checks equivalence with an SMT solver, reporting a counterexample when the
versions differ. Both give one verdict for the whole input space. Our Rule~3 is such a
check, but we only ask it on regions where the interval layer has already shown that both
versions return normally, so the rest of the domain still receives a verdict.

\paragraph{Differential symbolic execution and partitioning}
Differential symbolic execution~\citep{person2008dse} summarises both versions symbolically
and abstracts unchanged code with uninterpreted functions. ARDiff~\citep{badihi2020ardiff}
refines those uninterpreted functions on demand when the abstraction is too coarse to decide
equivalence. Our z3only baseline is plainer: whole-method SMT equivalence with no abstraction of common code. Partition-based regression
verification~\citep{boehme2013prv} is the closest in spirit: it splits the input space into
partitions with the same symbolic paths, and reports each partition as equivalent or
different, so the result is useful even when the whole program cannot be decided.
PASDA~\citep{glock2024pasda} extends partition-based differencing with best-effort
classifications for partitions that cannot be proved. We share the goal of per-region
verdicts, but our partitions are boxes produced by splitting at code constants, and many of
them are decided by an abstract interpretation without path enumeration, which also
decides loop pairs such as \code{loopThenZero}.

\paragraph{Abstract interpretation of differences}
\citet{partush2013abstract} analyse the \emph{correlating program} of two versions with
numerical abstract domains, so that the abstraction can express relations such as
$x_A = x_B$ directly. \citet{trostanetski2017modular} compute under- and
over-approximations of the input sets on which two versions differ, procedure by procedure,
and refine summaries on demand. Partush and Yahav can express relations between the versions inside the abstract domain,
which helps on loops. We deliberately analyse each
version separately with plain intervals and leave relations to the solver, at the price of
the \unk{} for \code{loopDiv/\_rewritten}.

\paragraph{Alternatives from the course}
A \emph{syntactic analysis} (a bytecode diff, or patterns for known refactorings) cannot
tell that deleting \code{if (d == 0) return 0} matters only at $(0,0)$; we use syntax only
to collect the split constants. A \emph{bounded static analysis} silently drops long traces,
which breaks the soundness that Rules~1 and~2 need, so we use the unbounded analysis with
widening. \emph{Concolic execution} and \emph{coverage-guided fuzzing} would find better
witnesses than random testing, but like differential
testing~\citep{mckeeman1998differential} they can refute equivalence and never prove a
region \pres{}. A \emph{relational domain} such
as octagons~\citep{mine2006octagon} on a product program could prove some preserved loop
pairs, but needs a product construction and much harder soundness and widening proofs. We
chose intervals, whose proofs we could complete, and moved the relational reasoning to the
solver.
